% %%%%%%%%%%%%%%%%%%%%%%%%%%%%%%%%%%%%%%%%%%%%%%%%%%%%%%%%%%%%
% % Task 02 : Single View Geometry and Distance Estimation
% %%%%%%%%%%%%%%%%%%%%%%%%%%%%%%%%%%%%%%%%%%%%%%%%%%%%%%%%%%%%

% \section{Overview}

% In this task, you will estimate real-world distances using a single
% image by -- well -- applying what you learnt in class. The task
% focuses on  linearity, cross-ratio, and the practical use
% of object localization (using a  deep learning tool) to infer
% distances from a single image. This topic is called single view metrology.

% \section{Background}

% You are provided with an image (see Fig.~\ref{fig:bicycleBBox}) of
% the way from the main building
% (MB) of IIT Bombay to the arch of knowledge (gyanam-paramam-dhyeyam) (``arch''). (As an aside, the arch encodes the motto of IITB).

% Take a moment to digest the provided image. Observe
% \begin{itemize} 
% \item the red line overlaid for reference, 
% \item the cross road (``intersection'') that has the palm tree on one side and a building
% on the other side
% \item a bicycle
% \end{itemize}

% Amongst the  objects, look for the street lights  towards the right. We refer to these as 
% ``poles'' in this assignment. The poles are numbered starting from the camera to the Arch. We will casually refer to the red line as equivalent to MB.
% We are
% going to assume that the cross road (aka intersection)  is a straight line passing
% through the base of the third pole parallel to the red
% line, and that the mb-to-arch line is perpendicular to
% the intersection line.


% %%%%%%%%%%%%%%%%%%%%%%%%%%%%%%%%%%%%%%%%%%%%%%%%%%%%%%%%%%%%
% %\section{Real life distances}
% %%The following real world distances are hypothesized:
% %\begin{itemize}
% %\item Distance from the red line to the arch: \textbf{160 metres}
% %\item Distance from  the red line to the intersection: \textbf{50 metres}
% %\end{itemize}


% %%%%%%%%%%%%%%%%%%%%%%%%%%%%%%%%%%%%%%%%%%%%%%%%%%%%%%%%%%%%


% %%%%%%%%%%%%%%%%%%%%%%%%%%%%%%%%%%%%%%%%%%%%%%%%%%%%%%%%%%%%
% \section{Who \href{https://en.wikipedia.org/wiki/Who_Moved_My_Cheese\%3F}{moved} my bicycle?}%

% \subsection{Task}
% Consider the bicycle in
% Fig.~\ref{fig:bicycleBBox}. Report the real-world distance between the
% bicycle and the arch.

% \begin{figure}[h!]
% \centering
% \includegraphics[width=0.45\linewidth]{bicycle_bb}
% \caption{A bounding box prediction from YOLOv8.}
% \label{fig:bicycleBBox}
% \end{figure}

% Use a pretrained deep learning model to detect the
% bicycle. This is an object detection task. Please refer to tutorial here, and stick to \textbf{YOLO} models: \url{https://docs.ultralytics.com/tasks/detect#models}
% Also look at the COCO dataset \url{https://cocodataset.org/#home} to find out the class identifiers for bicycles and persons.
% %The bounding box from the model will give you the coordinates of the bicycle. 

% \begin{studentSpace}
%     \vspace{10cm}
% \end{studentSpace}


% \subsection{Questions}

% You are expected to write answers inline in the student space provided. (see Section 5)

% \begin{enumerate}
% \item Use a tool from Google that tells you the real world distance from MB (the red line) to the arch.  Use a number in metres and approximate to the nearest distance divisible by 10. 
% \begin{studentSpace}
%     \vspace{3cm}
%     Measured \textbf{137.22 m} with Google Earth's measure tool, rounded to the nearest 10 - \textbf{ 140 m.}
% \end{studentSpace}

% \item Why can you NOT find the distance to the bicycle without making any assumption about the bicycle using the same tool? (Answer in one line.)
% \begin{studentSpace}
%     \vspace{3cm}
%     Google's ruler only works between things it can actually see and pin down on the satellite photo - buildings, roads, fixed spots. \textbf{The bicycle isn't on that map (it's a movable thing that wasn't there when the satellite photo was taken), so there's nothing to click on for it. To get its distance you first have to assume where it's standing} \end{studentSpace}

% \item Ensure that you draw a line in dark blue color (using an image annotation tool of your choice) that is used for obtaining the cross reference.  Show this line and the corresponding image in your submission. How did you draw the line? Keep in mind that cross-ratio is a \textbf{1D} concept and requires 4 collinear points. Make sure these points lie on the dark blue color line above.
% \begin{studentSpace}
%     \vspace{3cm}
%     First I found that vanishing point by extending the two edges of the path until they crossed. Then I drew the dark-blue line from that point straight down through the spot where the bicycle touches the ground. This line is parallel to the road in reality, so it meets the same vanishing point, and it passes right through the bicycle. This gave me my four points sitting on one straight line: the main building at the bottom A', then the bicycle B', then the arch C', then the vanishing point at the top V'. I drew it in MS Paint with the straight-line tool and read off each point's pixel position.
% \end{studentSpace}

% \item Ensure that you annotate 4 points of interest as A', B', C', V'  in sorted order of distance from the camera where B' corresponds to the biycle and V' corresponds to the vanishing point.  What do A' and B' correspond to?
% We recommend using MS Paint for annotation and identifying pixel coordinates.
% \begin{studentSpace}
%     \vspace{3cm}
%     \textbf{A′}is the main-building reference - the point where the dark-blue line starts. It's the anchor at 0 m. Paired with C′ (the arch at 140 m), A′C' gives you the two known real-world distances that make the whole cross-ratio solvable. A′ and C′ are the knowns. 

%     \textbf{B′} is the bicycle - the unknown we're solving for. B' corresponds to the specific pixel where the bicycle meets the ground.
% \end{studentSpace}

% \item Be sure to  explain in detail your method (and code) for finding the bicycle. Does the choice of reference point within the bounding box from the deep learning model affect distance estimation accuracy? Quantify.
% \end{enumerate}
% \begin{studentSpace}
%     \vspace{3cm}

%     \textbf{Finding the bicycle}

%     I used a pretrained YOLO detector (YOLO26 from Ultralytics), which already knows how to spot bicycles. My first run with the small model found nothing, because it shrinks the image to 640 px and the far-away bicycle disappeared. I switched to the large model (yolo26x) and a bigger 1280-px view, and it found the bicycle at 92% confidence.

%         \begin{verbatim}
%         from ultralytics import YOLO
%         model = YOLO("yolo26x.pt")
%         r = model("sample.jpg", imgsz=1280, conf=0.10)[0]
%         for b in r.boxes:
%             if int(b.cls) == 1:          # bicycle
%                 print(b.xyxy[0].tolist(), float(b.conf))
%         \end{verbatim}
        
%     I use the bottom-centre (the ground contact) as the bicycle's position.

%     \textbf{Does the reference point matter?}
    
%     Yes. Trying the top, middle, and bottom of the box gives 113.4 m, 115.4 m, and 117.1 m to the arch — a spread of about 3.8 m. Bottom-centre is correct (it's where the wheel meets the ground), so my answer is ≈117 m from the arch


% \end{studentSpace}


% %%%%%%%%%%%%%%%%%%%%%%
% \section{Your picture}

% In this section, take a picture of you and, alternately, your teammate
% (or some other person) similar to the one that has been provided in
% your locality. Identify two specific points of interest and tell us
% how far s/he is from the point of interest in the real world. We are
% expecting two pictures in your submission; one corresponding to the
% problem statement (often called as the first person view), and the
% second, a picture of you taking a picture of your teammate, i.e., a
% third person view. The second picture will show both the cameraman and
% the teammate (and ideally at the same time).

% \subsection{Tasks}

% \begin{enumerate}
% \item Capture and include the images in your submission.
% % \begin{studentSpace}
% %     \vspace{5cm}
% % \end{studentSpace}
% \item Detect the teammate using the same deep learning model.
% % \begin{studentSpace}
% %     \vspace{5cm}
% % \end{studentSpace}
% \item Estimate the distance between one of the objects of interest and
%   the team mate.
% \begin{studentSpace}
%     \vspace{5cm}
%     I captured a first-person view (the teammate standing in the room with the reference line) and a third-person view (a second photo showing me photographing him, both of us in frame, to document that the scene and the 120 cm were real). I detected the teammate with the same pretrained YOLO model used in the bicycle task (person = COCO class 0), which returned a bounding box at 0.94 confidence. I then estimated the distance from the teammate to the object of interest (a can/bottle placed on the floor) using the cross-ratio method.
% \end{studentSpace}
% \end{enumerate}

% You are expected to write answers inline.

% \subsection{Question}
% \begin{enumerate}

% \item As before, draw the line of interest (in blue color) marking all points of interest. 
% % \begin{studentSpace}
% %     \vspace{5cm}
% % \end{studentSpace}
  
% \item Explain how you got the real world distances.
% \begin{studentSpace}
% I placed four points on one floor tile seam: a reference mark $A'$, the object (a can) at $C'$, the teammate's foot contact at $B'$, and the vanishing point $V'$ where the floor lines converge. By distance from the camera: $A'$ (0 cm), $C'$ (60 cm), $B'$ (teammate, unknown), $V'$ (at infinity). I measured $A'$ to $C'$ on the floor with a tape: 60 cm. One known length suffices because the cross-ratio of four collinear points is preserved under projection, and with $V'$ at infinity it reduces to a simple ratio.

% Using $C(x_1,x_2,x_3,x_4)=\dfrac{|x_1x_2|\,|x_3x_4|}{|x_1x_3|\,|x_2x_4|}$ with points $A',C',B',V'$ and pixel positions $A'{=}3105$, $C'{=}2712$, $B'{=}2338$, $V'{=}1618$:
% \[ C=\frac{A'C'\cdot B'V'}{A'B'\cdot C'V'}=\frac{393\times720}{767\times1094}=0.337 \]
% In the real world ($V'$ at infinity, its terms cancel), with the teammate at distance $p$ from $A'$:
% \[ C=\frac{60\times\infty}{p\times\infty}=\frac{60}{p} \]
% Equating: $60/p=0.337 \Rightarrow p=178$ cm, so the teammate is $178-60=118$ cm behind the object. A tape measure gave 120 cm, so the estimate is accurate to within about 2 cm. The teammate was found with YOLO (person, COCO class 0), using the box bottom-centre as $B'$. The main error source is that contact pixel: the true contact gave 118 cm, the looser box bottom gave $\approx$100 cm --- amplified by the short 60 cm baseline.
% \end{studentSpace}

  
% \item Discuss practical challenges encountered while applying the projective
% geometry concepts to real-world images. Number all challenges and write independent challenges i.e. don't mix up the challenges.
% \begin{studentSpace}
%     \vspace{10cm}
    
% 1. Collinearity: My first arrangement had the teammate standing off the receding line, so the four points were not collinear; the cross-ratio is only valid for points on one straight line, so the line had to be redrawn through the foot contact toward the shared vanishing point.
% 2. In one attempt the teammate and the object sat at the same distance from the camera — only 3 pixels apart in depth — which blew the cross-ratio up and returned a meaningless negative distance. Separating them in depth fixed it.
% 3. Occlusion: When the teammate stood in front of the object, his body hid the object's floor contact, so it could not be read. I solved this by placing both reference points on the floor in front of the teammate.

% \end{studentSpace}


% \item What programs did you write and why?  
% \begin{studentSpace}
%     \vspace{5cm}
%     A YOLO detection script (detect.py, using Ultralytics YOLO26). Purpose: to locate the teammate automatically and read the bounding box / foot-contact pixel, rather than guessing it by eye — reproducible and consistent with the bicycle task. It runs the pretrained model on the photo and filters for the person class.
% \end{studentSpace}


% \end{enumerate}

% \section{Submit}

% Write a description of your method and all relevant aspects of your
% experiments. Write it in line in the text under the relevant section.  Your writing should be similar to  a blog post so that anyone in the class can fully understand the pain/joy/content you experienced (feel free to take copious input from an LLM but do not substitute an LLM for learning).

% The total length should be no less than 500 words and
% should normally exceed 800 words.  Your ultimate aim in writing your answers
% to mention important aspect of your assignment for the task, with
% emphasis on the last task.  This will include the methods employed
% listing all the software tools you employed, as well experimental
% aspects such as when you took the picture, how many pictures 
% you took, whether the method works in dawn, morning, and dusk, and so
% on. Be creative.  

% Do not submit any large file such as
% model weights and ancillary files that we can easily get from the
% internet to reproduce.

%%%%%%%%%%%%%%%%%%%%%%%%%%%%%%%%%%%%%%%%%%%%%%%%%%%%%%%%%%%%
% Task 02 : Single View Geometry and Distance Estimation
%%%%%%%%%%%%%%%%%%%%%%%%%%%%%%%%%%%%%%%%%%%%%%%%%%%%%%%%%%%%

%%%%%%%%%%%%%%%%%%%%%%%%%%%%%%%%%%%%%%%%%%%%%%%%%%%%%%%%%%%%
% HONOR CODE  (edit the bracketed parts; keep only if true)
%%%%%%%%%%%%%%%%%%%%%%%%%%%%%%%%%%%%%%%%%%%%%%%%%%%%%%%%%%%%
\section*{Honor Code}

I/we affirm that the work submitted here is our own and that we have not
engaged in any form of academic dishonesty.

\noindent\textbf{Sources used:}
\begin{itemize}
\item \textbf{LLM.} I used Anthropic's Claude as a tutor and debugging
  assistant throughout this task --- to understand the cross-ratio
  method, to fix the YOLO/Python environment, and to check the
  distance calculations. I did not substitute it for learning; all
  photographs, tape measurements, annotations, and final decisions are
  my own. Representative prompts: \emph{[e.g. ``explain how cross-ratio
  recovers a real distance from one image'', ``why does my YOLO run
  return no detections'', ``verify my cross-ratio calculation against
  these pixel coordinates'']}. Conversation link: \emph{[paste link]}.
\item \textbf{Software/tools.} Ultralytics YOLO26, Google Earth,
  MS~Paint, Python (NumPy), \texttt{pdflatex}.
\item \textbf{People.} Teammate: \emph{[name, if any]}.
\end{itemize}

%%%%%%%%%%%%%%%%%%%%%%%%%%%%%%%%%%%%%%%%%%%%%%%%%%%%%%%%%%%%

\section{Overview}

In this task, you will estimate real-world distances using a single
image by -- well -- applying what you learnt in class. The task
focuses on  linearity, cross-ratio, and the practical use
of object localization (using a  deep learning tool) to infer
distances from a single image. This topic is called single view metrology.

\section{Background}

You are provided with an image (see Fig.~\ref{fig:bicycleBBox}) of
the way from the main building
(MB) of IIT Bombay to the arch of knowledge (gyanam-paramam-dhyeyam) (``arch''). (As an aside, the arch encodes the motto of IITB).

Take a moment to digest the provided image. Observe
\begin{itemize} 
\item the red line overlaid for reference, 
\item the cross road (``intersection'') that has the palm tree on one side and a building
on the other side
\item a bicycle
\end{itemize}

Amongst the  objects, look for the street lights  towards the right. We refer to these as 
``poles'' in this assignment. The poles are numbered starting from the camera to the Arch. We will casually refer to the red line as equivalent to MB.
We are
going to assume that the cross road (aka intersection)  is a straight line passing
through the base of the third pole parallel to the red
line, and that the mb-to-arch line is perpendicular to
the intersection line.


%%%%%%%%%%%%%%%%%%%%%%%%%%%%%%%%%%%%%%%%%%%%%%%%%%%%%%%%%%%%
\section{Who \href{https://en.wikipedia.org/wiki/Who_Moved_My_Cheese\%3F}{moved} my bicycle?}%

\subsection{Task}
Consider the bicycle in
Fig.~\ref{fig:bicycleBBox}. Report the real-world distance between the
bicycle and the arch.

\begin{figure}[h!]
\centering
\includegraphics[width=0.45\linewidth]{bicycle_bb}
\caption{A bounding box prediction from YOLOv8.}
\label{fig:bicycleBBox}
\end{figure}

Use a pretrained deep learning model to detect the
bicycle. This is an object detection task. Please refer to tutorial here, and stick to \textbf{YOLO} models: \url{https://docs.ultralytics.com/tasks/detect#models}
Also look at the COCO dataset \url{https://cocodataset.org/#home} to find out the class identifiers for bicycles and persons.

\begin{studentSpace}
    \vspace{2cm}
\end{studentSpace}


\subsection{Questions}

You are expected to write answers inline in the student space provided. (see Section 5)

\begin{enumerate}
\item Use a tool from Google that tells you the real world distance from MB (the red line) to the arch.  Use a number in metres and approximate to the nearest distance divisible by 10. 
\begin{studentSpace}
    Measured \textbf{137.22 m} with Google Earth's measure tool, rounded to the nearest 10 --- \textbf{140 m}.
\end{studentSpace}

\item Why can you NOT find the distance to the bicycle without making any assumption about the bicycle using the same tool? (Answer in one line.)
\begin{studentSpace}
    Google's measure tool only works between things it can see and pin down on the satellite photo (buildings, roads, fixed spots). \textbf{The bicycle is not on that map (it is a movable object that was not there when the satellite imagery was taken), so there is nothing to click on for it; to get its distance you must first assume where it is standing.}
\end{studentSpace}

\item Ensure that you draw a line in dark blue color (using an image annotation tool of your choice) that is used for obtaining the cross reference.  Show this line and the corresponding image in your submission. How did you draw the line? Keep in mind that cross-ratio is a \textbf{1D} concept and requires 4 collinear points. Make sure these points lie on the dark blue color line above.
\begin{studentSpace}
    First I found the vanishing point by extending the two edges of the
    path until they crossed. Then I drew the dark-blue line from that
    point straight down through the spot where the bicycle touches the
    ground. This line is parallel to the road in reality, so it meets the
    same vanishing point, and it passes right through the bicycle. This
    gave me my four points on one straight line: the main building at the
    bottom $A'$, then the bicycle $B'$, then the arch $C'$, then the
    vanishing point at the top $V'$. I drew it in MS~Paint with the
    straight-line tool and read off each point's pixel position.

    \begin{center}
    % set this to your annotated (blue-line) bicycle image
    \includegraphics[width=0.45\linewidth]{bicycle_line}
    \end{center}
\end{studentSpace}

\item Ensure that you annotate 4 points of interest as A', B', C', V'  in sorted order of distance from the camera where B' corresponds to the biycle and V' corresponds to the vanishing point.  What do A' and B' correspond to?
We recommend using MS Paint for annotation and identifying pixel coordinates.
\begin{studentSpace}
    \textbf{$A'$} is the main-building reference --- the point where the
    dark-blue line starts. It is the anchor at 0~m. Paired with $C'$ (the
    arch at 140~m), $A'$ and $C'$ give the two known real-world distances
    that make the whole cross-ratio solvable; they are the knowns.

    \textbf{$B'$} is the bicycle --- the unknown we are solving for. $B'$
    corresponds to the specific pixel where the bicycle meets the ground
    (the bottom-centre of the bounding box).
\end{studentSpace}

\item Be sure to  explain in detail your method (and code) for finding the bicycle. Does the choice of reference point within the bounding box from the deep learning model affect distance estimation accuracy? Quantify.
\end{enumerate}
\begin{studentSpace}

    \textbf{Finding the bicycle.}
    I used a pretrained YOLO detector (YOLO26 from Ultralytics), which
    already recognises bicycles ($\texttt{bicycle}$ is class 1 in the COCO
    dataset). My first run with the small model found nothing, because it
    shrinks the image to 640~px and the far-away bicycle disappeared. I
    switched to the large model (\texttt{yolo26x}) and a larger 1280~px
    inference size, and lowered the confidence threshold; the bicycle was
    then detected at 92\% confidence. (The model also tagged the parked
    scooter as a motorcycle, confirming the bicycle I located is the
    distant one and not the scooter.)

\begin{verbatim}
from ultralytics import YOLO
model = YOLO("yolo26x.pt")
r = model("sample.jpg", imgsz=1280, conf=0.10)[0]
for b in r.boxes:
    if int(b.cls) == 1:          # COCO class 1 = bicycle
        print(b.xyxy[0].tolist(), float(b.conf))
\end{verbatim}

    Detected box: $(1111, 2088)$--$(1200, 2148)$. I use the bottom-centre
    (the ground contact) as the bicycle's position $B'$.

    \textbf{Does the reference point matter?}
    Yes. Trying the top, middle, and bottom of the box gives 113.4~m,
    115.4~m, and 117.1~m to the arch --- a spread of about 3.8~m.
    Bottom-centre is correct (it is where the wheel meets the ground), so
    my answer is $\approx$117~m from the arch (about 23~m from the main
    building). The estimate is even more sensitive to the placement of
    $C'$ and $V'$, because in this image the arch sits very close to the
    vanishing point.
\end{studentSpace}


%%%%%%%%%%%%%%%%%%%%%%
\section{Your picture}

In this section, take a picture of you and, alternately, your teammate
(or some other person) similar to the one that has been provided in
your locality. Identify two specific points of interest and tell us
how far s/he is from the point of interest in the real world. We are
expecting two pictures in your submission; one corresponding to the
problem statement (often called as the first person view), and the
second, a picture of you taking a picture of your teammate, i.e., a
third person view. The second picture will show both the cameraman and
the teammate (and ideally at the same time).

\subsection{Tasks}

\begin{enumerate}
\item Capture and include the images in your submission.
\item Detect the teammate using the same deep learning model.
\item Estimate the distance between one of the objects of interest and
  the team mate.
\begin{studentSpace}
    I captured a first-person view (the teammate standing in the room
    with the reference line, Fig.~\ref{fig:firstPerson}) and a
    third-person view (Fig.~\ref{fig:thirdPerson}, showing me
    photographing him with both of us in frame, to document that the
    scene and the 120~cm were real). I detected the teammate with the same
    pretrained YOLO model used in the bicycle task (person = COCO class 0),
    which returned a bounding box at 0.94 confidence. I then estimated the
    distance from the teammate to the object of interest (a can placed on
    the floor) using the cross-ratio method.

    \begin{center}
    \includegraphics[width=0.40\linewidth]{firstPersonView}
    \hfill
    \includegraphics[width=0.40\linewidth]{thirdPersonView}
    \end{center}
\end{studentSpace}
\end{enumerate}

You are expected to write answers inline.

\subsection{Question}
\begin{enumerate}

\item As before, draw the line of interest (in blue color) marking all points of interest. 
\begin{studentSpace}
    The dark-blue line runs along one floor tile seam through, in order of
    distance from the camera, $A'$ (a floor mark, 0~cm), $C'$ (the can,
    60~cm), $B'$ (the teammate's foot contact), and $V'$ (the vanishing
    point where the floor lines converge). All four points lie on this one
    straight line. See Fig.~\ref{fig:firstPerson}.

    \begin{figure}[h!]
    \centering
    % set these to your own image filenames
    \includegraphics[width=0.40\linewidth]{firstPersonView}
    \caption{First-person view with the blue line of interest and the
    four collinear points $A', C', B', V'$.}
    \label{fig:firstPerson}
    \end{figure}

    \begin{figure}[h!]
    \centering
    \includegraphics[width=0.40\linewidth]{thirdPersonView}
    \caption{Third-person view: the cameraman photographing the teammate.}
    \label{fig:thirdPerson}
    \end{figure}
\end{studentSpace}
  
\item Explain how you got the real world distances.
\begin{studentSpace}
    I placed four points on one floor tile seam: a reference mark $A'$,
    the object (a can) at $C'$, the teammate's foot contact at $B'$, and
    the vanishing point $V'$ where the floor lines converge. By distance
    from the camera: $A'$ (0~cm), $C'$ (60~cm), $B'$ (teammate, unknown),
    $V'$ (at infinity). I measured $A'$ to $C'$ on the floor with a tape:
    60~cm. One known length suffices because the cross-ratio of four
    collinear points is preserved under projection, and with $V'$ at
    infinity it reduces to a simple ratio.

    Using the cross ratio
    $C(x_1,x_2,x_3,x_4)=\dfrac{|x_1x_2|\,|x_3x_4|}{|x_1x_3|\,|x_2x_4|}$
    with the four points as $A', C', B', V'$ and pixel positions
    $A'=3105$, $C'=2712$, $B'=2338$, $V'=1618$:
    \[
        C=\frac{A'C' \cdot B'V'}{A'B' \cdot C'V'}
         =\frac{393\times 720}{767\times 1094}=0.337.
    \]
    In the real world ($V'$ at infinity, so its terms cancel), with the
    teammate's distance from $A'$ written as $p$:
    \[
        C=\frac{60\times \infty}{p\times \infty}=\frac{60}{p}.
    \]
    Equating image and world cross-ratios:
    $60/p = 0.337 \Rightarrow p = 178$~cm. The teammate therefore stands
    178~cm from $A'$, i.e.\ $178-60 = 118$~cm behind the object. I verified
    this with a tape measure: the true distance was 120~cm, so the
    single-view estimate was accurate to within about 2~cm.

    The teammate was located using the same pretrained YOLO model
    (person = COCO class 0), taking the bottom-centre of the detected box
    as the foot-contact point $B'$. The main source of error is the choice
    of that contact pixel: the true foot contact gave 118~cm, whereas the
    looser YOLO box bottom gave $\approx$100~cm --- a sensitivity made
    worse by the short 60~cm reference baseline.
\end{studentSpace}

\item Discuss practical challenges encountered while applying the projective
geometry concepts to real-world images. Number all challenges and write independent challenges i.e. don't mix up the challenges.
\begin{studentSpace}
\begin{enumerate}
\item \textbf{Collinearity.} My first arrangement had the teammate
  standing off the receding line, so the four points were not collinear;
  the cross-ratio is only valid for points on one straight line, so the
  line had to be redrawn through the foot contact toward the shared
  vanishing point.
\item \textbf{Equal-depth degeneracy.} In one attempt the teammate and
  the object sat at the same distance from the camera --- only 3 pixels
  apart in depth --- which blew the cross-ratio up and returned a
  meaningless negative distance. Separating them in depth fixed it.
\item \textbf{Occlusion.} When the teammate stood in front of the object,
  his body hid the object's floor contact, so it could not be read. I
  solved this by placing both reference points on the floor in front of
  the teammate.
\item \textbf{Reference-pixel ambiguity.} The bottom of the YOLO box did
  not coincide with the true foot contact; the box bottom gave
  $\approx$100~cm while the true contact gave 118~cm against a measured
  120~cm. Which pixel represents the subject matters more than the
  detection itself.
\item \textbf{Short reference baseline.} With $A'$ to $C'$ only 60~cm
  apart, every pixel of error was strongly amplified; a $\pm 20$~px shift
  in the foot point moved the result by about $\pm 10$~cm.
\item \textbf{Vanishing-point uncertainty.} The tiled floor had faint
  seams, making the vanishing point hard to locate precisely, and the
  cross-ratio is most sensitive to points near $V'$.
\end{enumerate}
\end{studentSpace}

\item What programs did you write and why?  
\begin{studentSpace}
\begin{enumerate}
\item \textbf{A YOLO detection script} (\texttt{detect.py}, Ultralytics
  YOLO26). Purpose: to locate the teammate (and earlier the bicycle)
  automatically and read the bounding box / foot-contact pixel, instead of
  guessing it by eye --- reproducible and consistent across both tasks. It
  runs the pretrained model on the photo and filters for the relevant
  class.
\item \textbf{A cross-ratio calculator} (Python, NumPy). Purpose: the
  cross-ratio arithmetic is error-prone by hand. The script takes the four
  pixel coordinates and the known 60~cm, computes the image cross-ratio,
  and solves for the teammate's distance. It also let me re-run the
  calculation while moving the foot-contact reference point, which is how I
  quantified the reference-pixel sensitivity reported above.
\end{enumerate}
\end{studentSpace}

\end{enumerate}

\section{Submit}

\begin{studentSpace}
\textbf{A short blog post on the experience.}

Single view metrology turns out to be one of those ideas that looks like
a magic trick and then becomes obvious once you see the machinery. The
whole task rests on a single fact from class: the cross-ratio of four
collinear points is invariant under projective transformation. A camera
is a projective transformation, so the cross-ratio I measure in pixels
equals the cross-ratio of the same four points in the real world. If
three of those points are known (or one known length plus the vanishing
point, which is the point at infinity), I can solve for the fourth. That
is the entire engine behind both parts of this task.

\textbf{The bicycle.} I started by measuring the main-building-to-arch
distance in Google Earth (137.22~m, rounded to 140~m). For the bicycle I
used a pretrained YOLO model rather than guessing its pixel location. My
first run quietly returned nothing: the small (nano) model downsizes the
image to 640~px, and the bicycle is so far down the path that it
effectively vanished. Switching to the large \texttt{yolo26x} model at
1280~px resolution and a lower confidence threshold recovered it at 92\%
confidence. I drew the dark-blue line through the bicycle's ground contact
toward the vanishing point, marked $A'$ (MB), $B'$ (bicycle), $C'$ (arch)
and $V'$, and the cross-ratio gave roughly 117~m from the bicycle to the
arch. The interesting wrinkle here is that the arch lies almost on top of
the vanishing point, which makes the estimate twitchy --- a few pixels of
error in $C'$ or $V'$ moves the answer by metres.

\textbf{My own picture (the part that taught me the most).} This is where
the theory met reality and lost a few rounds first. My initial photo had
my teammate standing beside the object rather than in front of it, so when
I computed the distance I got a meaningless negative number. The reason
was instructive: the teammate and the object were at essentially the same
depth from the camera (about three pixels apart along the receding line),
and the cross-ratio divides by that gap, so it exploded. The cure was to
arrange them at clearly different depths. The next problem was occlusion
--- when the teammate stood on the line in front of the sofa, his body hid
the object's contact with the floor, so I could not read its pixel. The
fix was to flip the layout: put both reference marks on the floor in front
of the teammate and stand him behind them, so nothing I needed was ever
blocked. With a can as the object, a tape-measured 60~cm between the two
floor marks, and the teammate detected by YOLO (person, COCO class 0), the
cross-ratio estimated 118~cm from the teammate to the can. I then measured
it with a tape: 120~cm. An error of about 2~cm from a single phone photo
felt like a genuine result.

\textbf{Tools and experiments.} I used Ultralytics YOLO26 for detection,
Google Earth for the campus distance, MS~Paint for annotation and reading
pixel coordinates, and a small Python/NumPy script for the cross-ratio.
The final indoor photo took several attempts --- the first few failed on
collinearity, equal depth, and occlusion before the geometry came out
clean. I shot in daytime indoor light. Lighting matters in two ways: low
or uneven light lowers YOLO's detection confidence, and it also washes out
the faint tile seams I rely on to find the vanishing point, so the same
method would be noticeably harder at dawn or dusk and most reliable in
bright, even daylight.

\textbf{What I take away.} The arithmetic was the easy part. The real work
was the geometry of the setup --- keeping the four points genuinely
collinear, separating the subject and object in depth, and avoiding
occlusion --- and realising that the single most error-sensitive decision
is which pixel I call the subject's ground contact. That choice moved my
indoor answer by nearly 20~cm, more than the detection model itself ever
did.
\end{studentSpace}

% Reminder: do NOT submit model weights (*.pt) or the torch/CUDA stack;
% they are easily reproducible from the internet.

